\section{The Ancient Argument: Thieves Do Not Undress a Corpse}

\subsection{Why the question arises inside Scripture}

The suggestion that the body was stolen is not a later sceptical invention that
the Fathers went out of their way to refute. It is in the Gospels. Matthew
reports that the chief priests asked Pilate for a guard because ``lest perhaps
his disciples come and steal him away and say to the people: He is risen from
the dead'' (Mt 27:64), and that after the event they paid the soldiers to say
``His disciples came by night and stole him away when we were asleep'' (Mt
28:13), adding that ``this word was spread abroad among the Jews even unto this
day'' (Mt 28:15). Whatever one makes of Matthew's account of the guard, the
evangelist is telling his readers that a theft explanation was in circulation
when he wrote.

An argument from the state of the burial cloths is therefore an argument
against a named and contemporary rival explanation, not a devotional
embellishment. That is how Chrysostom uses it.

\subsection{Chrysostom's argument, read at its own locus}

The eighty-fifth homily on John covers the crucifixion, the burial, and the
first ten verses of chapter 20. Having reached the burial, Chrysostom notes the
providential reasons for a nearby, new, sealed and guarded tomb, and closes
that catena of reasons with a sentence that sets up everything after it: ``Nor
was it for these reasons only that He was laid near, but also that the story
about the stealing might be proved false.'' The argument proper comes a page
later:

\begin{studyframe}\small
``When then she came and said these things, they hearing them, draw near with
great eagerness to the sepulchre, and see the linen clothes lying, which was a
sign of the Resurrection. For neither, if any persons had removed the body,
would they before doing so have stripped it; nor if any had stolen it, would
they have taken the trouble to remove the napkin, and roll it up, and lay it in
a place by itself; but how? they would have taken the body as it was. On this
account John tells us by anticipation that it was buried with much myrrh, which
glues linen to the body not less firmly than lead; in order that when thou
hearest that the napkins lay apart, thou mayest not endure those who say that
He was stolen. For a thief would not have been so foolish as to spend so much
trouble on a superfluous matter. For why should he undo the clothes? and how
could he have escaped detection if he had done so? since he would probably have
spent much time in so doing, and be found out by delaying and loitering. But
why do the clothes lie apart, while the napkin was wrapped together by itself?
That thou mayest learn that it was not the action of men in confusion or haste,
the placing some in one place, some in another, and the wrapping them together.
From this they believed in the
Resurrection.''\endnote{Chrysostom, \emph{Homilies on the Gospel of St.\
John}, Homily LXXXV.4, on Jn 19:38--20:10; NPNF, first series, vol.~14,
pp.~320--321, read in the printed volume. The homily's own scriptural heading
is Jn 19:16--18; the exposition runs to the end of the tomb scene.}
\end{studyframe}

This is worth taking apart, because it is a better argument than its
reputation.

\textbf{It has three distinct steps, not one.} The first is about motive: a
thief wants the body, and stripping it is work that does not serve his object.
The second is about physical difficulty, and it is the step usually dropped in
retellings---Chrysostom appeals to Jn 19:39--40, the hundred pounds' weight of
myrrh and aloes bound into the wrappings, and argues that this glues linen to
flesh ``not less firmly than lead.'' Unwrapping such a body is not a matter of
lifting a sheet. The third is about time and exposure: a man doing unnecessary
work in a guarded tomb increases his risk with every minute.

\textbf{It is an argument from the cost of an action, not from tidiness.}
Chrysostom does not say the tomb was neat. He says the cloths were in two
different states in two different places---``the placing some in one place,
some in another''---and that this pattern is inconsistent with men working in
confusion or haste. The inference is about the absence of hurry, and hurry is
what a thief in a sealed tomb must have.

\textbf{It is an argument from a state, not from a message.} Nothing in
Chrysostom suggests that the arrangement was a signal, a sign left
deliberately, or a communication of any kind. He is reasoning from physical
evidence to the behaviour that could and could not have produced it. This
distinction matters when the modern claim is examined in section~8, because
the modern claim is of an entirely different kind.

\subsection{What the argument establishes, and what it does not}

Taken at its own level, Chrysostom's argument is sound in structure and modest
in conclusion. It does not attempt to prove the Resurrection. It attempts to
eliminate one alternative account of an empty tomb, and it does so by asking
what a thief's behaviour would have looked like.

Three limits should be stated plainly.

It depends on the accuracy of the description in Jn 20:6--7, which no other
Gospel corroborates in its crucial detail. Luke's Peter sees cloths lying
alone; no Synoptic mentions the face-cloth.

It depends on the myrrh, and on a particular understanding of how the spices
were used. Jn 19:40 says the body was bound in the linen cloths ``with the
spices, as the manner of the Jews is to bury''; how a hundred pounds' weight of
myrrh and aloes was distributed is not stated, and Chrysostom's ``glues linen
to the body not less firmly than lead'' is a vivid inference rather than a
reported fact.

And it eliminates one explanation, not all of them. Removal of the body by
friends, by the tomb's owner, or by the authorities is not touched by the
motive argument in the same way, since such removal need not be hurried or
covert. Westcott, arguing from the same verse, felt the difference and wrote
that the deserted tomb ``bore the marks of perfect calm,'' concluding: ``It was
clear, therefore, that the body had not been stolen by enemies; it was scarcely
less clear that it had not been taken away by friends.''\endnote{Westcott,
\emph{The Gospel according to St.\ John}, vol.~2, p.~340, note on \emph{ch\=or.\
entetyl.\ eis hena top.}} The hedge in ``scarcely less clear'' is where the
argument's real strength ends.

\subsection{Westcott, Latham, and the modern form of the argument}

The nineteenth century produced a stronger version of the argument, and because
it is the ancestor of most modern English writing on the subject its origins
should be stated.

Its author was not a famous man. Henry Latham, Master of Trinity Hall,
Cambridge, published \emph{The Risen Master} in 1901 and credited the central
idea to a pamphlet by the Rev.\ Arthur Beard of St John's College. The
proposal is that the cloths were not merely undisturbed but \emph{collapsed}:

\begin{studyframe}\small
``In this recess, on the lower part of the ledge, lay the grave-clothes. They
were in no disorder, they were just as they were when Joseph and others had
wrapped them round the body of the Lord, only they were lying flat, fold over
fold, for the body was gone. On the raised part of the ledge at the far end,
all by itself, was the napkin that had gone round the head; this was not lying
flat, but was standing up a little, retaining the twirled form which had been
given it when it had been twined round the head of the
Lord.''\endnote{Latham, \emph{The Risen Master}, p.~34.}
\end{studyframe}

Bishop Westcott, whose commentary on John was published posthumously in 1908,
came to accept the substance of this. His son and editor prints two of his
father's marginal notes and a private letter, in which Westcott wrote that the
explanation of Jn 20:8 as the body's having ``passed through the heavy
wrappings as He later passed through doors'' was ``substantially correct'' but
added that he ``did not agree with his [Latham's] interpretation of
\emph{entetyligmenon}''; and, in a further marginal note: ``The undisturbed
grave-cloths show that the Lord had risen through and out of them. The face
cloth carefully rolled up, the action of the living
Lord.''\endnote{Arthur Westcott's editorial footnotes at B.~F. Westcott,
\emph{The Gospel according to St.\ John}, vol.~2, p.~339.}

Three observations, offered as this study's own assessment rather than as
anyone's finding.

The collapsed-cloths reading is an inference from silence about what the
disciples saw, dressed as a description. The text says the cloths were lying
and the face-cloth was apart and wrapped. It does not say the cloths retained a
body's outline, that they were flat, that they were undisturbed, or that no
spice had spilled. Latham is explicit that the absence of any mention of spices
is ``a significant feature in my case''---which is to say that his strongest
evidence is a thing the Gospels do not mention.

It is nevertheless not arbitrary, and it is the only reading that gives a
reason for the one contrast the verse actually draws. If the face-cloth
retained the shape a head had given it, then the reason it was not lying with
the linen cloths is that it had never been in the same place as the linen
cloths: it was where the head had been, and they were where the body had been.
On this reading \emph{eis hena topon} is doing real work, and the perfect
participle is doing exactly what perfects do.

And Latham himself---this is the point that matters most for section~8---
insisted that the participle rules out neat folding. He wrote it twice.

\subsection{The reception: how widely was this argument made?}

Chrysostom's is the fullest ancient statement of the anti-theft argument from
the cloths that this study located. The reception is treated in the next
section, including the material negative results: Fathers and medieval
commentators who had the occasion to make this argument and did not.
